\section{The state and a trace-preserving channel}\label{sec:applications}

We now combine the two independent inputs to the main state theorem: the
operational obstruction for the represented class and the product-free Choi
range of the explicit dimension-ten pair.

\begin{proof}[Proof of Theorem~\ref{thm:main-state}]
Take $\Phi_1,\Phi_2$ from \eqref{eq:phi-pair} and put
$Z=J(\Phi_2\circ\Phi_1)$.  Theorem~\ref{thm:main-pair} proves that $Z$ is
nonzero and that its range contains no nonzero product vector.  It is
positive because $\Phi_2\circ\Phi_1$ is CP.  Thus $\lambda=\tr Z>0$, and
$\rho=Z/\lambda$ is a density operator on $\C^{10}\otimes\C^{10}$ with the
same range as $Z$.  Lemma~\ref{lem:product-range} makes $\rho$ entangled.

Both maps are PPT.  Lemma~\ref{lem:choi-transfer}, with $a=b=c=10$, gives
\[
 \rho=(\Phi_1^\sharp\otimes\Phi_2)
       \left(\frac{\Omega_{10}\Omega_{10}^*}{\lambda}\right)
       \in\mathscr C.
\]
In particular, the representing vector is
$\Omega_{10}/\sqrt{\lambda}$; normalization is by the actual Choi trace.
Theorem~\ref{thm:zero-key} now gives \eqref{eq:main-gap} for every $n\geq1$
and every protocol and ideal bit in its stated scope, including the
purification and complete public record.  The same theorem gives
$K_D(\rho)=0$.  All matrices defining $\rho$ are specified by
\eqref{eq:explicit-matrices} and \eqref{eq:phi-pair}, so the state is
explicit.
\end{proof}

The maps $\Phi_1,\Phi_2$ need not preserve trace.  The next proposition
converts any square PPT pair into a trace-preserving PPT channel, while
retaining the pair's composite Choi matrix in a local corner.  The
two-block switch follows Christandl, Müller-Hermes, and Wolf
\cite[Section IV.A]{CMHW2019}.  The enlarged-input absorbing flag follows
Filippov's trace-preserving extension
\cite[Proposition 3, Equation (15)]{Filippov2021}.

\begin{proposition}\label{prop:channel-completion}
Let $\Psi_1,\Psi_2:\Mat m\to\Mat m$ be PPT maps.  There are explicitly
defined constants $c_1,c_2>0$, a trace-preserving PPT channel
$\Theta:\Mat{2m+1}\to\Mat{2m+1}$, and a real coordinate isometry
$W_0:\C^m\to\C^{2m+1}$ such that
\begin{equation}\label{eq:choi-corner}
 (W_0^*\otimes W_0^*)J(\Theta\circ\Theta)(W_0\otimes W_0)
       =c_1c_2J(\Psi_2\circ\Psi_1).
\end{equation}
If $J(\Psi_2\circ\Psi_1)$ is nonseparable, then $\Theta\circ\Theta$ is
not entanglement breaking.
\end{proposition}

\begin{proof}
Define
\begin{equation}\label{eq:channel-effects}
 A_i=\Psi_i^\dagger(I_m),\qquad
 c_i=\frac{1}{1+\tr A_i},\qquad
 F_i=I_m-c_iA_i \quad (i=1,2).
\end{equation}
The adjoint of a CP map is CP, so $A_i\succeq0$.  Since
$A_i\preceq(\tr A_i)I_m$, we have
$F_i\succeq(1-c_i\tr A_i)I_m=c_iI_m\succ0$.
The adjoint identity, using that $A_i$ is Hermitian, gives
\[
 \tr\Psi_i(B)=\tr(A_iB),\qquad
 \tr B-c_i\tr\Psi_i(B)=\tr(F_iB)
\]
for every complex matrix $B$.  Hence $c_i\Psi_i$ is trace nonincreasing
on positive inputs, and the positive effect $F_i$ records the missing
trace.

Use the standard-coordinate decomposition
\[
 \C^{2m+1}=\C^m_0\oplus\C^m_1\oplus\C e_\star,
\]
and let $W_0,W_1$ be the real coordinate inclusions.  For
$X\in\Mat{2m+1}$ write $X_{jj}=W_j^*XW_j$,
$x_\star=e_\star^*Xe_\star$, and $P_\star=e_\star e_\star^*$.  Set
\begin{align}\label{eq:channel}
 \Theta(X)={}&c_2W_0\Psi_2(X_{11})W_0^*
             +c_1W_1\Psi_1(X_{00})W_1^* \notag\\
 &+\bigl(\tr(F_1X_{00})+\tr(F_2X_{11})+x_\star\bigr)P_\star .
\end{align}
The formula discards the off-diagonal input blocks.  The two moving
branches are CP compositions.  For the flag branches, take spectral
decompositions $F_i=\sum_r\lambda_{ir}v_{ir}v_{ir}^*$ and set
\[
 K_{ir}=\sqrt{\lambda_{ir}}\,e_\star v_{ir}^*W_{i-1}^*
       \quad(i=1,2).
\]
Then
$\tr(F_iX_{i-1,i-1})P_\star=\sum_r K_{ir}XK_{ir}^*$.
The remaining flag term is $P_\star XP_\star$.  Thus $\Theta$ is CP.
Since $c_iA_i+F_i=I_m$, for every complex $X$,
\begin{align*}
 \tr\Theta(X)
 &=\tr\bigl[(c_1A_1+F_1)X_{00}\bigr]
   +\tr\bigl[(c_2A_2+F_2)X_{11}\bigr]+x_\star\\
 &=\tr X_{00}+\tr X_{11}+x_\star=\tr X.
\end{align*}
The discarded blocks have zero trace, so $\Theta$ is trace preserving.

Because $W_0,W_1,e_\star$ have real coordinates, output transposition in
\eqref{eq:channel} replaces $\Psi_i$ by $T_m\circ\Psi_i$ in the moving
branches and fixes the flag output.  Coordinate transpose is complex-linear:
it does not conjugate the possibly complex scalar multiplying $P_\star$.
The effects $F_i$ need not be real.  Each transposed moving branch is CP by
the PPT hypothesis, and the flag branches are CP as shown above.  Therefore
$T_{2m+1}\circ\Theta$ is CP, proving that $\Theta$ is PPT.

The flag is absorbing: $\Theta(P_\star)=P_\star$.  Applying
\eqref{eq:channel} twice to an input on the first block gives, for every
$B\in\Mat m$,
\begin{align*}
 \Theta^2(W_0BW_0^*)
 ={}&c_1c_2W_0\Psi_2(\Psi_1(B))W_0^*\\
 &+\bigl(\tr(F_1B)+c_1\tr(F_2\Psi_1(B))\bigr)P_\star .
\end{align*}
Compression to that block consequently gives
\begin{equation}\label{eq:square-corner}
 W_0^*\Theta^2(W_0BW_0^*)W_0
       =c_1c_2\Psi_2(\Psi_1(B)).
\end{equation}
Expanding in the coordinate matrix units of the input turns
\eqref{eq:square-corner} into \eqref{eq:choi-corner}.  The coordinate
qualification fixes the reference factor: for general complex isometries
$W,V$ and a compatible map $\Gamma$,
\[
 J\bigl(B\mapsto V^*\Gamma(WBW^*)V\bigr)
   =(W^{\T}\otimes V^*)J(\Gamma)(\overline W\otimes V).
\]
Here $W_0$ is real, so the reference factor is exactly the one displayed in
\eqref{eq:choi-corner}.

Local compression preserves separability, since it takes a positive product
$A\otimes B$ to
$(W_0^*AW_0)\otimes(W_0^*BW_0)$.  Thus separability of $J(\Theta^2)$
would imply separability of the positive multiple
$c_1c_2J(\Psi_2\circ\Psi_1)$ in \eqref{eq:choi-corner}.  This proves the
last assertion.
\end{proof}

\begin{proof}[Proof of Theorem~\ref{thm:main-channel}]
Apply Proposition~\ref{prop:channel-completion} with $m=10$ and
$\Psi_i=\Phi_i$ from \eqref{eq:phi-pair}.  Their composite Choi matrix is
nonseparable by Theorem~\ref{thm:main-pair} and
Lemma~\ref{lem:product-range}.  Formula \eqref{eq:channel} therefore gives
the required PPT channel on $\Mat{21}$.
\end{proof}

The product-free range assertion is used for the compressed matrix $Z$.
The full channel Choi matrix has a product contribution from the absorbing
flag: since $\Theta$ annihilates off-diagonal blocks involving $e_\star$ and
$\Theta(P_\star)=P_\star$, the matrix $P_\star\otimes P_\star$ is a direct
summand of $J(\Theta^2)$.  The channel conclusion is the nonseparability of
the full Choi matrix, certified by its entangled corner.
